\documentclass[a4paper]{article}
% Español (México) con XeLaTeX: fontspec para fuentes Unicode, polyglossia
% para la división silábica y los nombres del documento en la variante mexicana.
\usepackage{fontspec}
\usepackage{polyglossia}
\setdefaultlanguage[variant=mexican]{spanish}
% Los nombres chinos van como «pinyin (original)»: xeCJK da a los caracteres
% Han su propia fuente; sin ella XeLaTeX los omite con un aviso.
% FandolSong no cubre algunos caracteres de nombres (U+73FA); WenQuanYi Zen Hei sí.
\usepackage[AutoFallBack=true]{xeCJK}
\setCJKmainfont{FandolSong-Regular.otf}
\setCJKfallbackfamilyfont{\CJKrmdefault}{WenQuanYi Zen Hei}
% polyglossia no traduce los nombres de listings.
\AtBeginDocument{\providecommand{\lstlistingname}{}\renewcommand{\lstlistingname}{Listado}%
  \providecommand{\lstlistlistingname}{}\renewcommand{\lstlistlistingname}{Índice de listados}}
\usepackage{amsmath, amssymb}
\usepackage{graphicx}
\usepackage[margin=2.5cm]{geometry}
\usepackage[most]{tcolorbox}
\usepackage{etoolbox}
\usepackage{listings}
\usepackage{hyperref}
\usepackage{booktabs}
\usepackage{subcaption}
\usepackage{float}

\newtcolorbox{knowledgebox}[1]{enhanced, colback=blue!5!white, colframe=blue!75!black, colbacktitle=blue!75!black, coltitle=white, fonttitle=\bfseries, title={#1}, attach boxed title to top left={yshift=-2mm, xshift=2mm}, boxrule=1pt, sharp corners}
\newtcolorbox{importantbox}[1]{enhanced, colback=yellow!10!white, colframe=yellow!80!black, colbacktitle=yellow!80!black, coltitle=black, fonttitle=\bfseries, title={#1}, sharp corners}
\newtcolorbox{warningbox}[1]{enhanced, colback=red!5!white, colframe=red!75!black, colbacktitle=red!75!black, coltitle=white, fonttitle=\bfseries, title={#1}, sharp corners}

\lstset{language=Python, basicstyle=\ttfamily\small, keywordstyle=\color{blue}, stringstyle=\color{red!60!black}, commentstyle=\color{green!60!black}, breaklines=true, frame=single, numbers=left, numberstyle=\tiny\color{gray}, captionpos=b, extendedchars=true}

\newcommand{\notetitle}{[LLM Agents SP25] Multimodal Autonomous AI Agents}
\newcommand{\noteauthors}{Elaboradas a partir de la clase de Ruslan Salakhutdinov}
\newcommand{\notedate}{2026-04-02}
\newcommand{\videochannel}{Berkeley RDI}
\newcommand{\videopublishdate}{2025-03-11}
\newcommand{\videoduration}{1:17:42}
\newcommand{\videourl}{https://www.youtube.com/watch?v=RPINOYM12RU}
\newcommand{\videocoverpath}{cover.jpg}
\newcommand{\repourl}{https://github.com/hqhq1025/ai-course-notes}

\usepackage{fancyhdr}
\pagestyle{fancy}
\fancyhf{}
\fancyfoot[C]{\thepage}
\fancyfoot[R]{\footnotesize\color{gray}\href{\repourl}{AI Course Notes}}
\renewcommand{\headrulewidth}{0pt}
\renewcommand{\footrulewidth}{0pt}

\begin{document}
\begin{titlepage}
\centering
{\Large Notas del curso\par}\vspace{1.2cm}
{\huge\bfseries \notetitle\par}\vspace{0.8cm}
{\large \noteauthors\par}\vspace{0.3cm}
{\large \notedate\par}\vspace{1.2cm}
\ifdefempty{\videocoverpath}{}{\includegraphics[width=0.82\textwidth,height=0.45\textheight,keepaspectratio]{\videocoverpath}\par}
\vfill
\begin{tcolorbox}[width=0.9\textwidth, colback=black!2!white, colframe=black!60, sharp corners]
\textbf{Autor o canal del video}: \videochannel\par
\textbf{Fecha de publicación}: \videopublishdate\par
\textbf{Duración del video}: \videoduration\par
\ifdefempty{\videourl}{}{\textbf{Enlace del video}: \href{\videourl}{\nolinkurl{\videourl}}\par}
\textbf{Página del proyecto}: \href{\repourl}{\nolinkurl{\repourl}}\par
\end{tcolorbox}
\end{titlepage}
\tableofcontents
\newpage

\section{Ubicación del curso y pregunta central}

\subsection{Por qué importa esta clase}
Esta clase no busca presentar otro «modelo más grande», sino centrar la atención en un problema más de ingeniería: cuando el modelo entra a un entorno web real, cómo pasar de «saber responder preguntas» a «saber completar tareas». Ruslan resume el tema en tres puntos: evaluación, búsqueda y escalamiento de datos. Estos tres puntos son justamente donde la mayoría de los sistemas de Agentes encuentran primero obstáculos al llevarlos a producción.

\begin{importantbox}{Hilo conductor de la clase}
\begin{enumerate}
    \item \textbf{VisualWebArena}: cómo evaluar Agentes web multimodales en un entorno reproducible;
    \item \textbf{Inference-Time Search}: cómo lograr que el Agente sea mejor para probar, equivocarse y retroceder durante la inferencia;
    \item \textbf{Data Scaling}: cómo seguir mejorando la capacidad del Agente en condiciones de escasez de datos.
\end{enumerate}
\end{importantbox}

\begin{knowledgebox}{El cambio de paradigma de «conversar» a «ejecutar»}
Cuando la tarea pasa de responder preguntas a operar una página web, el modelo debe procesar al mismo tiempo la entrada visual, los cambios de estado, las restricciones de acción, los objetivos de largo plazo y la recuperación ante fallas. Este problema se acerca más a «Sequential Decision Making» que a la generación de texto en un solo turno.
\end{knowledgebox}

\begin{figure}[H]
\centering
\includegraphics[width=0.82\textwidth]{frame-01-visualwebarena.jpg}
\caption{Video aproximadamente en 00:06:50: el ponente presenta VisualWebArena y enfatiza la evaluación de Agentes multimodales en un entorno de interacción real}
\end{figure}

\subsection{Qué aporta esta clase directamente a quienes trabajan en la práctica}
El ponente insiste en un punto: el cuello de botella de los sistemas de Agentes no está solo en los parámetros del modelo, sino en la capacidad sistémica. Incluso si el modelo base es más fuerte, sin búsqueda, verificadores y datos de interacción de alta calidad, el límite superior de desempeño se alcanza muy pronto. A la inversa, aunque el modelo base no sea el más fuerte, con solo acertar en estas tres líneas es posible obtener beneficios considerables.

\begin{warningbox}{Error frecuente}
Equiparar directamente «la mejora de la capacidad del modelo» con «el aumento de la tasa de finalización de tareas» suele ser un error. En las tareas web, el espacio de estados es grande, los pasos son largos y el costo de los errores intermedios es alto, por lo que la mejora de la capacidad en un solo paso no se traduce automáticamente en una mayor tasa de éxito en cadenas largas.
\end{warningbox}

\subsection{Resumen de la sección}
Esta clase se centra en los tres pilares de la ingeniería de Agentes: evaluación, búsqueda y datos. No le interesa la calidad de una respuesta aislada, sino la tasa de éxito y la escalabilidad de las tareas de cadena larga.

\section{VisualWebArena: un banco de pruebas de evaluación orientado a tareas reales}

\subsection{De WebArena a VisualWebArena}
WebArena ya estructuró las tareas de interacción con páginas web, pero la representación puramente textual tiene limitaciones evidentes: gran parte de la semántica visual (el diseño, las imágenes, los gráficos, la posición de los botones) no se puede conservar por completo. VisualWebArena, en cambio, incorpora explícitamente la información visual y exige que el agente entienda a la vez las capturas de pantalla de la página y los elementos estructurados de la web.

\begin{importantbox}{Qué resuelve VisualWebArena}
\begin{itemize}
    \item las tareas ya no dependen solo del texto HTML, sino del contexto visual real;
    \item la evaluación ya no mide solo «si se puede escribir una respuesta textual», sino «si se puede completar una trayectoria de interacción»;
    \item el origen de los errores se acerca más al de los sistemas reales, e incluye clics equivocados, lecturas erróneas, detecciones omitidas y fallas de retroceso.
\end{itemize}
\end{importantbox}

\begin{knowledgebox}{Los tipos de tarea suelen incluir tres categorías}
\begin{enumerate}
    \item \textbf{Tipo de recuperación visual}: localizar el objetivo a partir de una imagen o un diseño;
    \item \textbf{Tipo de flujo entre sitios}: completar la tarea encadenando varias pestañas y varios sitios web;
    \item \textbf{Tipo de satisfacción de condiciones}: cumplir restricciones como precio, calificación o tiempo, y completar el envío.
\end{enumerate}
\end{knowledgebox}

\subsection{Por qué las tareas ``fáciles para los humanos'' siguen siendo difíciles para el agente}
El expositor mencionó que muchas tareas no son complejas para los humanos, pero siguen siendo difíciles para el agente; la razón principal es que ``la cadena de decisiones es demasiado larga y los errores se acumulan''. Un descuido aparentemente pequeño (hacer clic en el botón equivocado, leer mal un elemento de la página) puede desviar por completo la trayectoria posterior.

\begin{warningbox}{Las trampas de las métricas de evaluación}
Si solo se observa la tasa de éxito final, es fácil pasar por alto los modos de falla intermedios; si solo se observan los pasos intermedios, se puede sobrestimar la capacidad final de entrega. El método más confiable es una evaluación conjunta de ``métrica final + calidad de la trayectoria + tipo de error''.
\end{warningbox}

\begin{table}[H]
\centering
\small
\begin{tabular}{p{2.8cm}p{4.0cm}p{4.2cm}}
\toprule
\textbf{Dimensión de evaluación} & \textbf{Práctica común de un LLM de una sola ronda} & \textbf{Práctica al estilo VisualWebArena} \\
\midrule
Forma de entrada & prompt de texto puro & captura de pantalla de la página + estado estructurado de la web \\
Forma de salida & respuesta textual & secuencia de acciones de varios pasos \\
Criterio de corrección & coincidencia textual & resultado final de la tarea + consistencia de la trayectoria \\
Forma principal de falla & alucinación/omisión & operación errónea/falla de retroceso/desviación de la estrategia \\
\bottomrule
\end{tabular}
\caption{El cambio de paradigma de la evaluación de preguntas y respuestas textuales a la evaluación de tareas interactivas}
\end{table}

\subsection{Resumen de la sección}
El valor de VisualWebArena radica en transformar el objetivo de evaluación de ``saber decir'' a ``saber hacer'', e integrar en un marco unificado la percepción multimodal, la toma de decisiones de cadena larga y la estabilidad de ejecución.

\section{Formalización del entorno: POMDP, espacio de estados y espacio de acciones}

\subsection{Por qué formalizar}
El expositor modela al Web Agent como un POMDP, un paso muy importante. Esto nos permite describir el problema con un lenguaje unificado de aprendizaje por refuerzo y búsqueda, en lugar de tratar el comportamiento del sistema como un conjunto de scripts heurísticos.

\begin{figure}[H]
\centering
\includegraphics[width=0.82\textwidth]{frame-02-pomdp.jpg}
\caption{Video aproximadamente en 00:13:22: el expositor escribe el problema de interacción con la página web como states / actions / observations / transitions}
\end{figure}

\begin{knowledgebox}{Explicación intuitiva del POMDP en este contexto}
\begin{itemize}
    \item \textbf{State}: el estado real de la página web (incluida la información no visible);
    \item \textbf{Observation}: el contenido actualmente visible para el Agent (captura de pantalla, DOM, árbol de accesibilidad, etc.);
    \item \textbf{Action}: clic, escritura, desplazamiento, cambio de pestaña, detención, etc.;
    \item \textbf{Reward}: si la tarea se completó correctamente; suele ser escasa.
\end{itemize}
\end{knowledgebox}

\subsection{El diseño del espacio de acciones determina el límite superior del sistema}
Si el espacio de acciones es demasiado grueso, la capacidad expresiva resulta insuficiente; si es demasiado fino, el espacio de búsqueda explota. En la práctica suele necesitarse un diseño jerárquico de «acciones semánticas + acciones de bajo nivel»; por ejemplo, un subobjetivo como «buscar restaurante» lo entrega el planner, y el executor lo descompone en clics y escritura.

\begin{importantbox}{Principios de diseño del espacio de acciones}
\begin{enumerate}
    \item Las acciones atómicas deben ser ejecutables y verificables;
    \item Las acciones de alto nivel deben reducir la profundidad de búsqueda;
    \item La acción de detención debe ir acompañada de un criterio confiable para determinar la finalización;
    \item Cada acción debe llevar un registro rastreable, para facilitar la reproducción y la depuración.
\end{enumerate}
\end{importantbox}

\begin{table}[H]
\centering
\small
\begin{tabular}{p{3.2cm}p{4.1cm}p{3.7cm}}
\toprule
\textbf{Nivel de acción} & \textbf{Acciones típicas} & \textbf{Riesgo principal} \\
\midrule
Acciones de bajo nivel & click / type / scroll / hover & pasos largos; los errores locales se acumulan fácilmente \\
Acciones de nivel medio & open-site / query / filter & ambigüedad semántica; dificultad para vincular parámetros \\
Acciones de alto nivel & complete-subtask / stop & detención prematura o juicio erróneo de finalización \\
\bottomrule
\end{tabular}
\caption{Niveles de acción del Web Agent y sus riesgos}
\end{table}

\subsection{Retos de entrenamiento y evaluación bajo recompensas escasas}
En las tareas web, muchos reward solo se otorgan al final, por lo que la asignación de crédito es muy difícil. Si el sistema depende únicamente de la señal final, difícilmente aprende la capacidad de corrección fina en los pasos intermedios.

\begin{warningbox}{Efectos secundarios de la recompensa escasa}
\begin{itemize}
    \item El entrenamiento es inestable, con alta varianza en la actualización de la política;
    \item El modelo tiende a aprender shortcuts oportunistas;
    \item En las tareas de cadena larga, los errores intermedios son difíciles de penalizar a tiempo.
\end{itemize}
\end{warningbox}

\begin{figure}[H]
\centering
\includegraphics[width=0.82\textwidth]{frame-03-trajectory.jpg}
\caption{Video aproximadamente en 00:16:03: el expositor muestra la trayectoria de una tarea compleja y enfatiza la coherencia entre el razonamiento de cadena larga y la ejecución}
\end{figure}

\subsection{Resumen de la sección}
La formalización mediante POMDP hace que la tarea del Agent pase de ser «ingeniería de scripts» a ser un «problema de decisión». Pero la recompensa escasa y la explosión del espacio de acciones hacen que, con un solo paso de inferencia hacia adelante, sea difícil completar tareas complejas de manera estable.

\section{Búsqueda en el momento de la inferencia: de la ejecución de una sola trayectoria a la decisión con posibilidad de retroceso}

\subsection{Por qué es necesario introducir la búsqueda}
La conclusión que da el expositor es muy clara: sin capacidad de búsqueda, en cuanto el agente se desvía es muy difícil que regrese a la trayectoria correcta. Esto es especialmente cierto en escenarios web, donde una acción errónea cambia el estado de la página, el espacio de operaciones posteriores cambia en consecuencia, y el costo de recuperación es muy alto.

\begin{importantbox}{Beneficios directos de la búsqueda}
\begin{itemize}
    \item Permite evaluar en paralelo varias trayectorias candidatas;
    \item Permite el retroceso y reintentar por ramas;
    \item Convierte «acertar de un solo intento» en «elegir la mejor tras varios intentos».
\end{itemize}
\end{importantbox}

\subsection{Estrategia básica: Rejection Sampling}
El método más simple consiste en generar primero varias trayectorias y luego filtrarlas con un criterio de resultado final. Es sencillo de implementar y favorece la paralelización, pero no es muy eficiente, porque una gran cantidad de trayectorias inválidas desperdicia el presupuesto de inferencia.

\begin{knowledgebox}{Cuándo funciona Rejection Sampling}
Cuando la tarea es corta, el verificador es confiable y hay suficientes recursos para paralelizar, este método puede servir como referencia de bajo costo de ingeniería y aportar mejoras rápidas.
\end{knowledgebox}

\subsection{Estrategia avanzada: guía por valor y búsqueda en árbol}
Cuando aumenta la longitud de la tarea, el muestreo puro deja de funcionar rápidamente. En ese momento hay que introducir una función de valor o un Judge que estime el potencial de los estados intermedios, dando prioridad a expandir los nodos de mayor valor, para lograr «explorar mientras se podan ramas».

\begin{figure}[H]
\centering
\includegraphics[width=0.82\textwidth]{frame-07-search-eval.jpg}
\caption{Video en aproximadamente 00:34:39: el expositor explica cómo se evalúa el resultado de una trayectoria en una rama de búsqueda y se decide si continuar o terminar}
\end{figure}

\begin{table}[H]
\centering
\small
\begin{tabular}{p{2.8cm}p{3.8cm}p{4.4cm}}
\toprule
\textbf{Método} & \textbf{Ventajas} & \textbf{Desventajas} \\
\midrule
Ejecución de una sola trayectoria & Baja latencia, implementación sencilla & Los errores no son recuperables, la tasa de éxito depende mucho del ruido de una sola decisión \\
Rejection Sampling & Favorece la paralelización, fácil de poner en producción & Desperdicio de cómputo notable, la ganancia disminuye en tareas largas \\
Búsqueda guiada por valor & Alta eficiencia de muestreo, permite retroceso & Ingeniería compleja, sensible a la calidad del Judge \\
\bottomrule
\end{tabular}
\caption{Comparación de estrategias de razonamiento del agente}
\end{table}

\begin{warningbox}{La búsqueda no es gratuita}
La búsqueda aumenta de manera notable el costo de inferencia. Si no hay un buen mecanismo de poda y reutilización de nodos, el sistema puede caer en un estado de «presupuesto agotado sin haber convergido».
\end{warningbox}

\subsection{Resumen de la sección}
En las tareas web de cadena larga, la capacidad de búsqueda no es opcional, sino la fuente central de la robustez. Lo importante no es «si se busca», sino «cómo se busca de manera eficiente».
\section{Desacoplamiento entre el modelo de planeación y el modelo de ejecución}

\subsection{Por qué se necesita la división Planner-Executor}
En la segunda mitad, el expositor destacó un enfoque jerárquico: el modelo de nivel superior se encarga de planear, y el modelo de nivel inferior se encarga de ejecutar. De esta manera, una tarea compleja puede dividirse en subtareas localmente controlables, lo que reduce el ruido de tomar decisiones directas paso a paso.

\begin{figure}[H]
\centering
\includegraphics[width=0.82\textwidth]{frame-06-planning.jpg}
\caption{Video, aproximadamente 00:27:40: el expositor explica que primero el planning model produce un plan y luego el modelo de ejecución lo concreta en acciones}
\end{figure}

\begin{knowledgebox}{Beneficios del desacoplamiento}
\begin{itemize}
    \item La capa de planeación puede imponer restricciones sobre el objetivo global;
    \item La capa de ejecución puede realizar acciones localmente óptimas para la página actual;
    \item La localización de errores es más clara, lo que facilita la depuración y la reversión.
\end{itemize}
\end{knowledgebox}

\subsection{Cómo diseñar la interfaz jerárquica}
Una interfaz aplicable en la práctica suele contener tres partes: la descripción del subobjetivo, los límites de ejecución y el criterio de finalización. Sin una interfaz clara, la capa de planeación y la capa de ejecución terminan interfiriéndose entre sí, lo que produce el fenómeno de ``el plan está bien escrito pero la ejecución se desvía''.

\begin{importantbox}{Campos mínimos de la interfaz jerárquica}
\begin{enumerate}
    \item \textbf{subgoal}: el objetivo de la etapa actual;
    \item \textbf{allowed actions}: el conjunto de acciones permitidas;
    \item \textbf{stop condition}: el criterio para determinar si se completó o falló;
    \item \textbf{fallback rule}: la condición para regresar a la capa superior en caso de anomalía.
\end{enumerate}
\end{importantbox}

\subsection{Cuándo conviene volver a planear}
En cuanto la ejecución dispare una anomalía de alta confianza (el elemento no existe, un salto de página anómalo, exceso del presupuesto), el sistema debe regresar a tiempo a la capa de planeación, en lugar de continuar ciegamente con la siguiente acción.

\begin{warningbox}{Fallas comunes de la capa de planeación}
\begin{itemize}
    \item La granularidad del plan es demasiado gruesa y la capa de ejecución no puede aplicarlo;
    \item La granularidad del plan es demasiado fina y se pierde el valor de la jerarquía;
    \item Falta la regla de ``cuándo volver a planear'', lo que provoca que el error siga propagándose.
\end{itemize}
\end{warningbox}

\subsection{Resumen de la sección}
El desacoplamiento Planner-Executor es un patrón de ingeniería clave para los agentes de cadena larga. Lo que realmente determina el resultado es el diseño de la interfaz y la estrategia de replaneación, y no solo cambiar a un modelo más grande.

\section{Verifier y LLM-as-a-Judge: incorporar la evaluación al razonamiento}

\subsection{De la evaluación al final a la evaluación intermedia}
Si el éxito solo se determina al final de la trayectoria, la retroalimentación llega demasiado tarde. El ponente propone adelantar el verifier a los pasos intermedios, de modo que el sistema pueda juzgar, durante el razonamiento mismo, si «vale la pena continuar por el camino actual».

\begin{importantbox}{Las tres funciones del verifier}
\begin{enumerate}
    \item Puntuación de rutas: sirve de base para podar la búsqueda;
    \item Detección de errores: identifica estados que se han desviado de forma evidente;
    \item Filtrado de datos: selecciona trayectorias de alta calidad utilizables para el entrenamiento.
\end{enumerate}
\end{importantbox}

\subsection{El valor práctico de LLM-as-a-Judge}
Hacer que el LLM evalúe la trayectoria suele ser más estable que hacer que genere directamente la trayectoria óptima. Esto se debe a que «juzgar si algo es correcto» suele ser más fácil que «acertar de un solo paso», una lección práctica que este apartado reitera en varias ocasiones.

\begin{knowledgebox}{Elementos clave del prompt del juez}
\begin{itemize}
    \item El objetivo y las restricciones de la tarea;
    \item La trayectoria completa (o un resumen de los pasos clave);
    \item Una plantilla de salida explícita: aprobado o rechazado, junto con la justificación;
    \item Una estimación de confianza, que la estrategia de búsqueda pueda usar.
\end{itemize}
\end{knowledgebox}

\subsection{Confiabilidad y calibración del juez}
El juez no constituye una verdad absoluta. Si el criterio de evaluación es ambiguo o el prompt se ha desviado de su propósito original, el juez puede llegar a una conclusión errónea con alta confianza y, con ello, desviar la búsqueda.

\begin{warningbox}{Tres riesgos del juez}
\begin{itemize}
    \item Manipulación de la recompensa: el modelo aprende a complacer el estilo del juez en lugar de completar la tarea de verdad;
    \item Sesgo de dominio: el juez pierde precisión ante plantillas de páginas web con las que no está familiarizado;
    \item Desajuste de confianza: una puntuación alta no implica una tasa de éxito alta.
\end{itemize}
\end{warningbox}

\begin{figure}[H]
\centering
\includegraphics[width=0.82\textwidth]{frame-08-full-pipeline.jpg}
\caption{Minuto aproximado 01:12:10 del video: la planificación, la ejecución y el verificador forman un ciclo cerrado completo; la evaluación ya no ocurre únicamente al final de la tarea}
\end{figure}

\subsection{Resumen de la sección}
El verifier hace que el agente pase de «actuar primero y juzgar después» a «juzgar mientras actúa». Sin embargo, la confiabilidad del juez requiere una calibración continua, pues de lo contrario puede orientar al sistema hacia una dirección de optimización equivocada.
\section{Ampliación de datos: de trayectorias escasas a un ciclo de entrenamiento sostenible}

\subsection{Por qué los datos son el cuello de botella}
Los datos de interacción de alta calidad para un Web Agent son más costosos que los corpus de texto puro. Cada trayectoria contiene el estado del entorno, la secuencia de acciones, los resultados de las herramientas y el veredicto final, por lo que tanto la recolección como la depuración tienen un costo alto.

\begin{importantbox}{Dificultades de los datos}
\begin{itemize}
    \item La interacción real es lenta, por lo que el costo por muestra es alto;
    \item las trayectorias son largas, lo que hace complejas la anotación y la verificación;
    \item el entorno cambia con el tiempo, por lo que los datos antiguos se vuelven obsoletos.
\end{itemize}
\end{importantbox}

\subsection{Recolección y filtrado en línea}
El expositor recomienda muestrear en línea en un entorno real o en una simulación de alta fidelidad, y luego aplicar un filtrado por capas mediante un verifier o un Judge: primero se descartan los fracasos evidentes, después se conservan las trayectorias exitosas de alta confianza y, por último, se muestrea de forma equilibrada según el tipo de tarea.

\begin{knowledgebox}{Un pipeline de datos ejecutable}
\begin{enumerate}
    \item Muestreo: generación paralela de trayectorias con múltiples políticas;
    \item filtrado: veredicto del resultado final más veredicto del proceso;
    \item reponderación: se aumenta el peso de las muestras difíciles;
    \item entrenamiento: actualización mediante SFT o RL;
    \item repetición: pruebas A/B en línea para verificar si de verdad hay mejora.
\end{enumerate}
\end{knowledgebox}

\subsection{Estrategia combinada en las etapas de entrenamiento}
En la práctica, la combinación común es «base con SFT + refuerzo mediante búsqueda + ajuste fino con RL». El SFT ofrece un punto de partida estable, la búsqueda aporta trayectorias de exploración de alta calidad y el RL refina la política con la retroalimentación de la interacción.

\begin{table}[H]
\centering
\small
\begin{tabular}{p{2.4cm}p{3.9cm}p{4.7cm}}
\toprule
\textbf{Etapa} & \textbf{Objetivo} & \textbf{Consideraciones} \\
\midrule
SFT & Aprender la sintaxis básica de las acciones y el formato de la tarea & Evitar que solo memorice plantillas e ignore la generalización en cadenas largas \\
Search-augmented & Mejorar la calidad de la exploración y la proporción de trayectorias exitosas & Controlar el presupuesto de inferencia y evitar el muestreo excesivo \\
RL / actualización de la política & Alinear la tasa de éxito final y la robustez & Evitar el reward hacking y el desplazamiento de la distribución \\
\bottomrule
\end{tabular}
\caption{Rutas de combinación de entrenamiento comunes en un Web Agent}
\end{table}

\begin{warningbox}{Riesgo implícito en la ampliación de datos}
Si solo se persigue la tasa de éxito, el sistema puede sesgarse hacia las «tareas fáciles», lo que provoca que el desempeño se deteriore en los escenarios difíciles. La ampliación de datos debe hacerse junto con la gestión de la distribución de dificultad de las tareas.
\end{warningbox}

\subsection{Resumen de la sección}
El núcleo de la ampliación de la capacidad del Agent no es solo la mejora del modelo, sino un ciclo de datos sostenible. El muestreo en línea, un filtrado confiable y el equilibrio de dificultad determinan si el sistema puede avanzar de manera estable.

\section{Interpretación de los resultados: mejora del desempeño, brecha entre modelos y costo computacional}

\subsection{De una línea base baja a un nivel cercano al humano}
El expositor mostró que, mediante mejores estrategias de razonamiento y mecanismos de evaluación, el desempeño del sistema puede mejorar de manera notable a partir de un nivel bajo, y en algunas configuraciones acercarse al desempeño humano. Este resultado muestra que el «diseño del sistema» es igual de decisivo para la mejora del desempeño.

\begin{figure}[H]
\centering
\includegraphics[width=0.82\textwidth]{frame-04-human-level.jpg}
\caption{Video aproximadamente en 00:22:23: el expositor muestra la mejora del desempeño y menciona el rango cercano al human level}
\end{figure}

\begin{importantbox}{El verdadero significado de estos resultados}
No significa que el Agente «ya haya resuelto la automatización web», sino que muestra que, dentro del marco correcto de evaluación y búsqueda, las tareas de cadena larga tienen un margen claro de optimización.
\end{importantbox}

\subsection{Evolución entre generaciones de modelos}
El expositor mencionó la tendencia de mejora de Gemini Pro, GPT-4 y modelos posteriores en este tipo de tareas. La capacidad de los modelos en efecto ha avanzado, pero el sistema de apoyo de «planear-ejecutar-verificar» sigue siendo una condición necesaria.

\begin{figure}[H]
\centering
\includegraphics[width=0.82\textwidth]{frame-05-model-comparison.jpg}
\caption{Video aproximadamente en 00:25:15: el expositor discute la mejora en el desempeño de las tareas entre Gemini Pro, GPT-4 y modelos posteriores}
\end{figure}

\begin{knowledgebox}{Cómo interpretar la mejora de los modelos}
\begin{itemize}
    \item Los modelos más fuertes reducen la tasa de error por paso;
    \item pero la estabilidad de la cadena larga sigue dependiendo del mecanismo de búsqueda y retroceso;
    \item mientras más fuerte es el modelo, más se necesita una evaluación y una gobernanza más estrictas, para evitar indicadores infladamente altos.
\end{itemize}
\end{knowledgebox}

\subsection{Equilibrio entre costo y beneficio}
La búsqueda y la verificación elevan la tasa de éxito, pero también aumentan el costo de tokens y de llamadas al entorno. Un sistema en producción debe gestionar de forma explícita el equilibrio entre «el beneficio de la tasa de éxito y el gasto en costo».

\begin{warningbox}{Errores comunes en la gestión de costos}
\begin{itemize}
    \item Perseguir solo la tasa de éxito, sin considerar el costo por tarea;
    \item pasar por alto el gasto marginal de los reintentos por fallo;
    \item no contar con una estrategia de clasificación de tareas, de modo que incluso las tareas simples pasan por una ruta de búsqueda intensiva.
\end{itemize}
\end{warningbox}

\subsection{Resumen de la sección}
Los resultados experimentales muestran que el desempeño del Agente puede mejorar de forma notable mediante un método sistemático. Pero pasar de la investigación a la puesta en producción todavía requiere optimizar a la vez la tasa de éxito, la latencia y el costo, en lugar de perseguir un solo indicador.
\section{Modos de falla y líneas de defensa de ingeniería}

\subsection{Falla típica: imposibilidad de recuperarse tras una acción errónea}
El ponente puso, durante la sesión de preguntas y respuestas, un ejemplo muy representativo: cuando el modelo ejecuta una acción errónea y carece de búsqueda y retroceso, resulta muy difícil autocorregirse después, y termina por desviarse por completo del objetivo de la tarea.

\begin{figure}[H]
\centering
\includegraphics[width=0.82\textwidth]{frame-09-failure-mode.jpg}
\caption{Video en aproximadamente 01:17:20: el ponente discute la dificultad de recuperarse tras una acción errónea y enfatiza la necesidad de mecanismos de búsqueda y retroceso}
\end{figure}

\begin{importantbox}{El ciclo mínimo de recuperación ante fallas}
\begin{enumerate}
    \item Detección: identificar rápidamente que la trayectoria actual se desvió;
    \item Reversión: volver al estado confiable más reciente;
    \item Replanificación: actualizar los subobjetivos y las restricciones de acción;
    \item Revisión: determinar mediante el verifier si se volvió a la trayectoria correcta.
\end{enumerate}
\end{importantbox}

\subsection{Clasificación de los modos de falla}
\begin{table}[H]
\centering
\small
\begin{tabular}{p{2.7cm}p{3.8cm}p{4.2cm}}
\toprule
\textbf{Tipo de falla} & \textbf{Síntomas típicos} & \textbf{Medidas de corrección de ingeniería} \\
\midrule
Error de percepción & Lectura equivocada de botones, campos de formulario o contenido de imágenes & Verificación multimodal, umbral de confianza de elementos \\
Error de planificación & Orden equivocado de subobjetivos u omisión de restricciones & Plantillas de plan, verificador de subobjetivos \\
Error de ejecución & Clics, entradas o detenciones equivocadas de manera continua & Registro de repetición, instantáneas de estado y reversión \\
Error de evaluación & El Judge da un veredicto equivocado de éxito o fracaso & Cruce entre varios Judges, calibración mediante revisión manual por muestreo \\
\bottomrule
\end{tabular}
\caption{Modos de falla comunes del Web Agent y estrategias de corrección}
\end{table}

\begin{warningbox}{No trates las ``fallas'' como muestras de ruido}
Las trayectorias de falla son una de las fuentes de datos más valiosas. El sistema debe archivar de manera sistemática los tipos de falla y realimentarlos al flujo de entrenamiento y evaluación, formando un ciclo de mejora continua.
\end{warningbox}

\subsection{La colaboración humana sigue siendo necesaria}
Para tareas de alto riesgo, el agente debe tener la capacidad de ``solicitar aclaración'' y de ``escalar a una persona''. La persona no debe aparecer solo como el último respaldo, sino intervenir de manera oportuna cuando aumenta la incertidumbre.

\begin{knowledgebox}{Ejemplos de condiciones que activan la colaboración humano-máquina}
\begin{itemize}
    \item Dos discrepancias consecutivas de alta confianza (el planner y el verifier no coinciden);
    \item Un campo clave no se puede determinar (relacionado con pagos, identidad o privacidad);
    \item Se rebasa el presupuesto pero la tarea no converge.
\end{itemize}
\end{knowledgebox}

\subsection{Resumen de la sección}
Los modos de falla del agente pueden gestionarse de manera sistemática. Un sistema verdaderamente confiable no es aquel que ``nunca se equivoca'', sino uno cuyos errores son ``detectables, recuperables y aprovechables para el aprendizaje''.

\section{Frontera de investigación: de agentes web a agentes ejecutores multimodales generales}

\subsection{Abstracciones de acción de nivel más alto}
El ponente propone una dirección que merece una inversión a largo plazo: diseñar para los agentes abstracciones semánticas de nivel más alto, de modo que el modelo no tenga que buscar desde acciones atómicas de click/type en cada tarea. Esto puede reducir de manera considerable la profundidad de búsqueda y la complejidad de ingeniería.

\begin{importantbox}{Beneficios potenciales de las abstracciones de nivel alto}
\begin{itemize}
    \item Reducen el ruido de las acciones de bajo nivel;
    \item Mejoran la generalización entre sitios;
    \item Crean una interfaz unificada para la colaboración entre múltiples agentes.
\end{itemize}
\end{importantbox}

\subsection{Búsqueda en paralelo y ejecución en múltiples instancias}
El ponente también mencionó la importancia de la paralelización: la misma tarea puede explorarse en múltiples instancias siguiendo rutas distintas, para después devolver la trayectoria óptima al flujo principal. Esto coincide con la tendencia de crecimiento del ``test-time compute'' en la era de los modelos grandes.

\begin{knowledgebox}{Puntos clave del diseño de la ejecución en paralelo}
\begin{enumerate}
    \item Las subtareas pueden reproducirse de forma independiente;
    \item Se comparten señales de evaluación intermedias;
    \item Se unifica la estrategia de resolución de conflictos (por ejemplo, tomar la ruta de mayor valor).
\end{enumerate}
\end{knowledgebox}

\subsection{El siguiente paso de los benchmarks de evaluación}
La próxima generación de benchmarks necesita cubrir de manera más realista: sitios web dinámicos, páginas personalizadas, sesiones de larga duración, interacción entre dispositivos y colaboración con múltiples roles. De lo contrario, el sistema mejorará en benchmarks estáticos pero fallará en el entorno de producción real.

\begin{warningbox}{La brecha estructural entre los benchmarks y el despliegue real}
Los productos web cambian de manera continua, los elementos se cargan de forma dinámica y las restricciones de permisos y sesión son complejas; cualquier benchmark fijo puede quedar obsoleto con rapidez. El sistema de evaluación debe admitir actualización continua y seguimiento de versiones.
\end{warningbox}

\subsection{Resumen de la sección}
El punto de avance a largo plazo de los agentes multimodales está en la abstracción de acciones, la búsqueda en paralelo y el sistema de evaluación dinámica. La competencia futura ocurrirá cada vez más en el nivel de la ingeniería de sistemas, y no solo en el nivel de la escala de parámetros.
\section{Análisis de caso: de la descripción de la tarea a la trayectoria ejecutable}

\subsection{El ejemplo de ``reservar un restaurante que cumpla las restricciones''}
La clase recurre varias veces a tareas de búsqueda de restaurantes para ilustrar las dificultades de ejecución del agente. Este tipo de tareas parece simple, pero en páginas web reales suele incluir restricciones múltiples: ubicación, precio, calificación, horario de atención, preferencia de cocina y el proceso de pedido. Si el sistema no cuenta con una gestión explícita del estado, es fácil que olvide condiciones a la mitad del proceso.

\begin{importantbox}{Estructura recomendada para descomponer la tarea}
\begin{enumerate}
    \item \textbf{Extracción de restricciones}: convertir los requisitos en lenguaje natural en campos estructurados;
    \item \textbf{Recuperación de candidatos}: generar un conjunto de candidatos dentro del sitio o entre sitios;
    \item \textbf{Verificación de evidencia}: verificar una por una las restricciones duras (como el tope de precio);
    \item \textbf{Ejecución de la acción}: entrar al proceso de pedido y completar el llenado de los campos clave;
    \item \textbf{Confirmación final}: confirmar por partida doble, mediante el estado de la página y el verifier, que la tarea se completó.
\end{enumerate}
\end{importantbox}

\begin{knowledgebox}{Por qué son clave las ``restricciones estructuradas''}
Si las restricciones no se estructuran, el agente tiende a perder información clave a lo largo de una cadena larga y solo conserva preferencias vagas. Al almacenar las restricciones explícitamente como pares clave-valor, puede hacerse una verificación de consistencia antes de cada paso de ejecución, lo que reduce de forma notable el costo de tener que rehacer trabajo por ``descubrir hasta el final que no se cumplía la condición''.
\end{knowledgebox}

\subsection{Control de calidad a nivel de trayectoria}
En esta tarea, que un paso individual sea correcto no significa que la trayectoria sea correcta. Por ejemplo, el modelo puede elegir primero un restaurante que se ve bien, pero cuyo precio u horario disponible no cumple los requisitos. El control de calidad a nivel de trayectoria debe verificar la ``consistencia con el objetivo'', y no solo la ``ejecutabilidad de la acción''.

\begin{table}[H]
\centering
\small
\begin{tabular}{p{2.8cm}p{3.5cm}p{4.4cm}}
\toprule
\textbf{Capa de verificación} & \textbf{Objeto de verificación} & \textbf{Errores comunes} \\
\midrule
Consistencia de campos & Precio, horario, ubicación, calificación & Campos omitidos, confusión de unidades, restricciones en conflicto \\
Consistencia de página & Si la página actual contiene la evidencia buscada & Error al identificar el elemento, lectura incorrecta de la información de la tarjeta \\
Consistencia del proceso & Si se entró a la ruta correcta de pedido & Sitio equivocado, entrada equivocada, detención prematura \\
Consistencia final & Si el resultado de la tarea coincide con el requisito original & Éxito parcial pero incumplimiento global \\
\bottomrule
\end{tabular}
\caption{Control de calidad a nivel de trayectoria en la tarea de búsqueda de restaurantes}
\end{table}

\begin{figure}[H]
\centering
\includegraphics[width=0.82\textwidth]{frame-07-search-eval.jpg}
\caption{Video aproximadamente en 00:34:39: comparar entre trayectorias candidatas antes de seguir expandiendo, lo que refleja la estrategia de ``evaluar antes de ejecutar''}
\end{figure}

\subsection{Plantillas de ejecución reutilizables}
Para reducir la necesidad de explorar desde cero en cada tarea, puede mantenerse una ``biblioteca de plantillas de tareas''. La plantilla no es un guion fijo, sino un ``esqueleto de proceso que puede restringirse'': define el orden de las fases, los puntos de verificación clave y la estrategia de reversión. Así se conserva la flexibilidad y, a la vez, se reduce de forma notable la varianza en tareas largas.

\begin{warningbox}{Los límites de usar plantillas}
Una plantilla demasiado rígida perjudica la generalización, y una demasiado laxa no logra restringir los errores. Lo razonable es usar la plantilla en los pasos de alto riesgo (pago, identidad, página de confirmación) y dejar la exploración para los pasos de bajo riesgo (recuperación y ordenamiento de candidatos).
\end{warningbox}

\subsection{Resumen de la sección}
Lo esencial en las tareas web complejas no es ``lograr que el modelo adivine mejor'', sino ``lograr que la trayectoria sea más controlable''. Las restricciones estructuradas, la verificación a nivel de trayectoria y las plantillas reutilizables son los tres trabajos básicos para convertir un sistema de demostración en un sistema estable.
\section{Lista de verificación de ingeniería: convertir un prototipo de investigación en un sistema de producción}

\subsection{Diez preguntas que deben responderse antes del lanzamiento}
En la etapa de prototipo de investigación, el sistema suele buscar mostrar casos de éxito lo antes posible; pero al pasar a producción, es necesario anteponer la robustez, el costo y la capacidad de auditoría. La siguiente lista puede usarse como plantilla de revisión previa al lanzamiento.

\begin{table}[H]
\centering
\small
\begin{tabular}{p{0.8cm}p{3.9cm}p{5.8cm}}
\toprule
\textbf{\#} & \textbf{Pregunta} & \textbf{Ejemplo de criterio de aprobación} \\
\midrule
1 & ¿Los límites de la tarea son claros? & Las restricciones de entrada, las acciones prohibidas y las condiciones de terminación están todas definidas en la documentación \\
2 & ¿El estado puede reproducirse? & Cada observation/action/result de cada paso puede reproducirse y auditarse \\
3 & ¿El presupuesto de búsqueda es controlable? & Se configuran un número máximo de pasos, una ramificación máxima y una estrategia de reversión por tiempo de espera agotado \\
4 & ¿El Judge está calibrado? & La precisión alcanza el objetivo en el conjunto de control fuera de línea, y existe un mecanismo de revisión periódica \\
5 & ¿La falla puede recuperarse automáticamente? & Admite una recuperación de tres etapas: detección de anomalías, reversión y replanificación \\
6 & ¿El disparador de intervención humana está definido con claridad? & Existe una estrategia de escalamiento obligatorio para escenarios de pago, privacidad y baja confianza \\
7 & ¿El retorno de datos forma un ciclo cerrado? & Tanto las trayectorias exitosas como las fallidas pueden retornar y participar en la siguiente ronda de entrenamiento \\
8 & ¿El costo puede explicarse? & Cada tipo de tarea cuenta con estadísticas de costo de tokens, latencia y llamadas externas \\
9 & ¿La evaluación se actualiza de manera continua? & Tanto el benchmark como el tráfico real se monitorean de manera continua \\
10 & ¿La seguridad y el cumplimiento están cubiertos? & Las operaciones sensibles cuentan con privilegio mínimo y registros de auditoría \\
\bottomrule
\end{tabular}
\caption{Lista mínima de ingeniería antes de la conversión en producto de un Web Agent}
\end{table}

\subsection{Los tres tipos de costos de ingeniería que más se subestiman}
\begin{enumerate}
    \item \textbf{Costo de gestión de estado}: un prototipo sin estado es rápido, pero la complejidad aumenta abruptamente al incorporar reversión y estado con múltiples ramificaciones;
    \item \textbf{Costo de evaluación}: construir un conjunto de control de alta calidad y calibrar el Judge de manera continua toma más tiempo que entrenar una demo;
    \item \textbf{Costo de operación}: la estructura de las páginas web cambia con frecuencia, por lo que las reglas de análisis y el mapeo de acciones requieren mantenimiento continuo.
\end{enumerate}

\begin{knowledgebox}{Por qué la ``operación'' se convierte en el costo principal}
Muchos equipos, en la etapa de PoC, solo observan el desempeño del modelo e ignoran la falla de las reglas provocada por los cambios en las páginas web. Lo que realmente consume mano de obra después del lanzamiento suele ser problemas de operación continua, como cambios de selector, rediseños del flujo de la página y actualizaciones de los mecanismos anti-automatización.
\end{knowledgebox}

\subsection{De impulsado por métricas a impulsado por valor}
El ponente mencionó ``economically valuable tasks'', lo que nos recuerda que no basta con mirar la puntuación del benchmark, sino que hay que incorporar el valor de la tarea al objetivo de optimización. Una tarea con puntuación alta pero valor bajo no debería ocupar un presupuesto de inferencia considerable; por el contrario, una tarea de alto valor debería permitir un mayor costo de búsqueda y una colaboración humana más intensa.

\begin{importantbox}{Recomendación de programación escalonada por valor}
\begin{itemize}
    \item Tareas de bajo valor: estrategia ligera + presupuesto bajo + falla rápida;
    \item Tareas de valor medio: búsqueda moderada + reversión automática;
    \item Tareas de alto valor: búsqueda intensiva + verificación múltiple + intervención humana.
\end{itemize}
\end{importantbox}

\subsection{Resumen de la sección}
La dificultad central de la conversión en producto de un Agent no está en ``si puede ejecutar una demo'', sino en ``si puede operar de manera continua, estable, controlable y auditable''. La lista de ingeniería y la programación escalonada por valor son las palancas clave para superar este umbral.

\section{Síntesis y ampliación}

\subsection{Conclusiones clave}
\begin{importantbox}{Conclusiones aplicables de esta clase}
\begin{enumerate}
    \item Primero construir una evaluación de interacción real, y después hablar de la optimización del modelo;
    \item la búsqueda en tiempo de inferencia es el elemento central de la robustez en tareas de cadena larga;
    \item Verifier/Judge es el componente clave que conecta la búsqueda, el entrenamiento y el filtrado de datos;
    \item el ciclo cerrado de datos determina si la capacidad puede seguir creciendo;
    \item los mecanismos de recuperación ante fallos y la colaboración humano-máquina deben incorporarse desde etapas tempranas del diseño del sistema.
\end{enumerate}
\end{importantbox}

\subsection{Tabla de síntesis}
\begin{table}[H]
\centering
\small
\begin{tabular}{p{2.6cm}p{3.9cm}p{4.3cm}}
\toprule
\textbf{Tema} & \textbf{Puntos clave del ponente} & \textbf{Acción de ingeniería} \\
\midrule
Evaluación & VisualWebArena mide la capacidad de interacción real & Construir una métrica conjunta de resultado final + trayectoria + tipo de error \\
Búsqueda & La ejecución de una sola trayectoria es inestable; se necesitan decisiones que puedan retroceder & Poner en producción búsqueda jerárquica con control de presupuesto \\
Verificador & Judge es mejor para ``evaluar'' que para ``hacer'' & Adelantar el verifier a los pasos intermedios de la inferencia \\
Ampliación de datos & El muestreo en línea más el filtrado forman un ciclo cerrado & Construir el retorno de ejemplos fallidos y el reponderado de dificultad \\
Robustez & La recuperación ante fallos y la colaboración humana no pueden omitirse & Agregar mecanismos de reversión, replanificación y activación de aclaraciones \\
\bottomrule
\end{tabular}
\caption{Resumen de la aplicación en ingeniería de la clase Multimodal Autonomous AI Agents}
\end{table}

\subsection{Resumen de la sección}
Esta clase no ofrece un solo algoritmo, sino una metodología de ingeniería para agentes: usar la evaluación real como ancla, usar la búsqueda y el verificador para mejorar la robustez de la inferencia, y usar el ciclo cerrado de datos para mantener la velocidad de iteración a largo plazo.

\subsection{Lecturas adicionales}
\begin{itemize}
    \item Artículos y código de referencia relacionados con VisualWebArena / WebArena;
    \item investigación sobre LLM-as-a-Judge y verificadores a nivel de trayectoria;
    \item comparación de métodos de Search-Augmented Inference en tareas de agentes;
    \item prácticas de recolección de datos de interacción en línea y filtrado de trayectorias de alta confianza;
    \item marcos de planificación multiagente y ejecución paralela (para tareas largas de alta complejidad).
\end{itemize}

\end{document}
